\documentclass[a4paper,12pt]{ctexart}
\usepackage{xcolor}
\usepackage[top=1.8cm, bottom=1.8cm, left=1.8cm, right=1.8cm]{geometry}
\usepackage{array}
\usepackage{longtable}
\usepackage{hyperref}
\hypersetup{colorlinks=true,linkcolor=blue!70!black}
\linespread{1.35}

\definecolor{donegreen}{RGB}{40,140,80}
\definecolor{physics}{RGB}{180,60,30}
\definecolor{chemistry}{RGB}{30,100,160}
\definecolor{biology}{RGB}{20,130,60}
\definecolor{earth}{RGB}{140,90,40}

\newcommand{\tagphysics}{\textcolor{physics}{\small [物理]}}
\newcommand{\tagchemistry}{\textcolor{chemistry}{\small [化学]}}
\newcommand{\tagbiology}{\textcolor{biology}{\small [生物]}}
\newcommand{\tagearth}{\textcolor{earth}{\small [地科]}}
\newcommand{\done}{\textcolor{donegreen}{\textbf{【已写】}}}

\title{\textcolor{blue!70!black}{\Huge\textbf{自然科学小故事 · 课程大纲}}}
\author{}
\date{}

\begin{document}
\maketitle
\thispagestyle{empty}

\section*{\textcolor{blue!70!black}{设计思路}}

本模块面向北京地区四年级起点的孩子，用讲故事的口吻介绍自然科学。每课围绕一个孩子生活中能见到、能问到、能动手试的现象展开，从具体体验过渡到背后的科学道理。不做公式推导，不要求背诵，只求唤起好奇心。

选课的三个原则：
\begin{enumerate}
  \item \textbf{从身边出发}——孩子见过、摸过、经历过
  \item \textbf{可以追问"为什么"}——每个现象背后有一条清晰的科学线索
  \item \textbf{学科均衡}——物理、化学、生物、地球科学各有覆盖
\end{enumerate}

\vspace{0.3cm}
\textcolor{gray}{\small 当前规划28课，可根据实际教学进度增删。}

\newpage

\section*{\textcolor{orange!80!black}{一、物理（力学与运动）}}

\begin{enumerate}

\item[\textbf{第1课}] \textbf{牛顿与苹果——万有引力} \done \\
\tagphysics \\
\textbf{内容：} 牛顿在苹果树下的思考，引出"地球在拉万物"的引力概念。扩展到月亮为什么不掉下来、宇航员为什么在空间站里"飘"。延伸到火箭如何利用反冲运动对抗引力飞向太空（第一/第二宇宙速度、嫦娥探月）。\\
\textbf{推荐理由：} 引力是孩子最早能感受到的"看不见的力"——跳起来被拉回来、东西往下掉。苹果故事有画面感，航天扩展激发对太空的想象。

\item[\textbf{第2课}] \textbf{阿基米德的浴缸——浮力的秘密} \done \\
\tagphysics \\
\textbf{内容：} 古希腊国王的金冠疑案，阿基米德洗澡时的"尤里卡"瞬间。引出浮力原理：排开多少水，就受多少托力。延伸到铁船为什么不沉（空心造型排开大量水）、飞机如何靠机翼升力"浮"在空中（伯努利原理）。\\
\textbf{推荐理由：} 每个孩子都玩过水、折过纸船，对"什么浮什么沉"有天然好奇。从浴缸到万吨巨轮，跨度震撼。

\item[\textbf{第3课}] \textbf{瓦特与烧水壶——蒸汽的力量} \done \\
\tagphysics \\
\textbf{内容：} 小瓦特在厨房看烧水壶盖子被蒸汽顶得"噗噗"跳。长大后改良蒸汽机，推动工业革命。延伸：今天的核电站——核裂变烧水→蒸汽→汽轮机→发电，原理和瓦特的烧水壶一脉相承。\\
\textbf{推荐理由：} 烧水是每个孩子都能观察到的日常场景，连接到一个改变世界的发明。再延伸到核电站，展现"大道理藏在小现象里"的科学之美。

\end{enumerate}

\section*{\textcolor{orange!80!black}{二、生物（生命的奇妙）}}

\begin{enumerate}

\item[\textbf{第4课}] \textbf{小蝌蚪找妈妈——生命的变形记} \done \\
\tagbiology \\
\textbf{内容：} 从经典童话切入，讲述青蛙的四阶段变态发育：卵→蝌蚪→长腿→成蛙。引出两栖动物的概念和特征（幼体用鳃、成体用肺+皮肤）。\\
\textbf{推荐理由：} 借孩子熟悉的故事降低认知门槛。变态发育是大自然最直观的"形态变化"，冲击力强、好记。

\item[\textbf{第5课}] \textbf{种子是怎么长大的？——植物生长的秘密} \\
\tagbiology \\
\textbf{内容：} 一颗豆子埋进土里，浇水、晒太阳，几天后冒出了小芽——它是怎么"知道"往上长的？引出种子结构（胚芽、胚根、子叶）、发芽条件（水、空气、温度）、向光性。延伸到参天大树都是从一颗种子开始的。\\
\textbf{推荐理由：} 学校科学课通常有种豆子实验，孩子有直接经验。从"看不见的土里发生什么"激发探究欲。

\item[\textbf{第6课}] \textbf{为什么跑步后会喘气？——呼吸的秘密} \\
\tagbiology \\
\textbf{内容：} 体育课刚跑完的感受——心跳变快、喘粗气。引出细胞需要氧气来"燃烧"食物产生能量（呼吸作用），就像汽车需要汽油+氧气。延伸到为什么运动员肺活量大、高原上为什么容易喘。\\
\textbf{推荐理由：} 每个孩子都经历过，但没想过"为什么"。从自身身体感受切入，是最直接的生物学入口。

\item[\textbf{第7课}] \textbf{吃下去的饭去哪了？——消化系统旅行记} \\
\tagbiology \\
\textbf{内容：} 一口饭从嘴到胃到小肠到大肠的"奇幻漂流"。口腔里唾液已经开始消化、胃酸有多强、小肠绒毛像毛巾一样吸收营养。用旅行故事的方式串起消化系统。\\
\textbf{推荐理由：} "肚子里发生了什么"是每个孩子都好奇过的问题。用故事化叙事把抽象器官变成旅行站点。

\item[\textbf{第8课}] \textbf{秋天树叶为什么变黄变红？} \\
\tagbiology \\
\textbf{内容：} 北京的孩子每年都经历——银杏金黄、香山红叶。春夏叶子绿是因为叶绿素在"上班"（光合作用），秋天日照变短、气温下降，叶绿素分解，被掩盖的类胡萝卜素（黄/橙）和花青素（红/紫）显露出来。\\
\textbf{推荐理由：} 北京秋天是一流教材。孩子亲眼看到颜色变化，只需要有人告诉他"为什么"。

\item[\textbf{第9课}] \textbf{小猫小狗为什么要睡觉？——动物的休息与修复} \\
\tagbiology \\
\textbf{内容：} 家里的猫一天睡十几个小时，狗也动不动就趴着。睡眠不是偷懒——大脑在清理代谢废物、身体在修复细胞。不同动物的睡眠策略：长颈鹿每天只睡两小时，蝙蝠倒挂着睡大半天。延伸到人为什么要早睡。\\
\textbf{推荐理由：} 孩子对宠物有感情，从他们关心的动物切入讲生物节律，接受度极高。

\item[\textbf{第10课}] \textbf{人为什么会发烧？——身体里的"战争"} \\
\tagbiology \\
\textbf{内容：} 发烧不是"病"，是免疫系统在战斗——大脑主动调高体温来抑制病菌繁殖、加速白细胞活动。引出免疫系统的基本概念：白细胞是"战士"，抗体是"通缉令"，发烧是"升高战场温度"。延伸到退烧药的原理（帮助身体暂时降低体温设定点）、什么时候需要吃退烧药。\\
\textbf{推荐理由：} 每个孩子都发过烧，但从没想过发烧是"自己人干的好事"。把生病变成故事，既能消除恐惧，又能理解身体的智慧。

\end{enumerate}

\section*{\textcolor{orange!80!black}{三、化学（生活中的物质变化）}}

\begin{enumerate}

\item[\textbf{第11课}] \textbf{为什么肥皂能洗手？——水的"皮肤"与肥皂分子} \done \\
\tagchemistry \\
\textbf{内容：} 水的表面张力（水面像一层橡皮膜）、水黾能在水面行走。肥皂分子的"双面结构"——亲水端和亲油端——如何破坏表面张力、包裹油污和细菌。延伸到洗手为什么比光用水冲有效90\%以上。\\
\textbf{推荐理由：} 洗手是每天必做的事，但从没想过肥皂和水之间发生了什么。化学从最日常的动作开始。

\item[\textbf{第12课}] \textbf{面包为什么又软又蓬松？——酵母的"呼吸"} \\
\tagchemistry \\
\textbf{内容：} 咬一口馒头或面包，里面全是小窟窿——这些窟窿是酵母"吹"出来的！酵母吃糖，放出二氧化碳气体，气泡被面筋困在面团里，加热后膨胀定型。延伸到小苏打发面的原理、为什么蒸出来的和烤出来的口感不同。\\
\textbf{推荐理由：} 孩子天天吃面食，但从面团到馒头的神奇变化他们很少想过。还可以在家和爸妈一起做一次馒头，变成亲子实验。

\item[\textbf{第13课}] \textbf{可乐里的气泡从哪来？——二氧化碳与压力} \\
\tagchemistry \\
\textbf{内容：} 拧开可乐瓶盖"嘶——"一声，气泡涌上来。为什么没开瓶时看不到气泡？因为高压把CO₂"逼"进了水里（溶解度）。开瓶后压力降低，气体跑出来了。延伸到汽水是怎么做出来的、为什么摇晃后再开会喷出来。\\
\textbf{推荐理由：} 可乐几乎每个孩子都喝过，气泡是直观的"气体在水里"的现象。成本低，现场演示效果好。

\item[\textbf{第14课}] \textbf{切开的苹果为什么变黑？——氧化反应} \\
\tagchemistry \\
\textbf{内容：} 刚切好的苹果，过了一会儿变黄变黑，看起来不好吃了。这是苹果里的酚类物质和空气中的氧气发生了氧化反应（酶促褐变）。滴几滴柠檬汁就能延缓变黑——维生素C是抗氧化剂。延伸到铁生锈也是氧化、人的衰老也和氧化有关。\\
\textbf{推荐理由：} 每天吃水果都可能遇到，是孩子在餐桌上就能观察到的化学现象。柠檬汁"救"苹果是一个可以当场做的小实验。

\item[\textbf{第15课}] \textbf{盐是从哪儿来的？——蒸发、结晶与人体之需} \\
\tagchemistry \\
\textbf{内容：} 盐不只是调料——人体没有盐神经就无法传递信号、肌肉就不能收缩。盐从哪来？三种来源：海水晒盐（把海水围在盐田里让太阳晒干→水蒸发→盐结晶）、井盐（古代四川用钻井取地下卤水熬盐）、岩盐（亿万年前海洋干涸留下的盐层，可以挖矿开采）。延伸到为什么运动后要喝淡盐水、动物也会找盐吃（盐渍地是野生动物的"餐厅"）。\\
\textbf{推荐理由：} 天天吃却不知来处的东西，最值得追问。结晶实验可以当堂演示（盐水蒸发后盐粒重现），非常直观。

\end{enumerate}

\section*{\textcolor{orange!80!black}{四、物理（光、热、声、电）}}

\begin{enumerate}

\item[\textbf{第16课}] \textbf{镜子里的世界——光的反射} \\
\tagphysics \\
\textbf{内容：} 为什么镜子里的"你"左右是反的？为什么哈哈镜里人会变胖变瘦？引出光的反射定律：入射角等于反射角。平面镜、凹面镜、凸面镜的不同效果。延伸到汽车后视镜是凸面镜（看更广）、牙医的小镜子、潜望镜原理。\\
\textbf{推荐理由：} 孩子天天照镜子，但"为什么左右反了上下没反"是一个看似简单却值得深究的问题。镜子是光学的最佳日常入口。

\item[\textbf{第17课}] \textbf{影子为什么跟着我走？——光沿直线传播} \\
\tagphysics \\
\textbf{内容：} 光走直线，被身体挡住就形成影子。早上影子长、中午影子短——因为太阳的角度变化。延伸到日晷（古代利用影子计时的"太阳钟"）、皮影戏（中国非遗，用影子讲故事）、日食和月食（天体级别的"影子游戏"——月亮挡住太阳就是日食，地球挡住月亮就是月食）。\\
\textbf{推荐理由：} 影子是孩子最早的"科学玩具"——踩影子、手影游戏。从玩乐到天文，一条直线串起来。

\item[\textbf{第18课}] \textbf{彩虹是从哪儿来的？——光的折射与色散} \\
\tagphysics \\
\textbf{内容：} 雨后天空的彩虹——为什么总是弯的？为什么有七种颜色？太阳光看起来是白的，其实是七色光的"合体"。雨后空气中有小水滴，光穿过水滴时被折射+反射+再折射，不同颜色折射角度不同，被"拆开"了。延伸到三棱镜实验、CD背面也能看到彩色。\\
\textbf{推荐理由：} 彩虹是孩子最早接触的"光学魔术"。近乎童话的外观下，藏着光的色散原理。可以配合三棱镜做课堂演示。

\item[\textbf{第19课}] \textbf{温度计为什么能测温度？——热胀冷缩} \\
\tagphysics \\
\textbf{内容：} 发烧时妈妈把温度计夹在你腋下——里面的红色液柱为什么自己升上去？引出液体热胀冷缩：温度升高，液体分子运动加快、间距增大、体积膨胀。酒精温度计和水银温度计的区别。延伸到夏天铁轨会变长需要留缝隙、高压电线夏天变松冬天变紧。\\
\textbf{推荐理由：} 生病必接触温度计，但从没想过"它怎么知道我在发烧"。热胀冷缩还能解释很多生活中的"奇怪事"。

\item[\textbf{第20课}] \textbf{为什么先看见闪电后听见打雷？——光与声的速度赛跑} \\
\tagphysics \\
\textbf{内容：} 每个雷雨天都会注意到这个现象。光速约每秒30万公里（一秒钟能绕地球七圈半），声速约每秒340米。闪电和雷声同时发生，光几乎瞬间到达，声音却要"跑"好一阵。可以教孩子自己估算雷的距离：看到闪电后开始数秒，每数3秒约等于1公里远。延伸到超音速飞机（突破音障时的音爆）。\\
\textbf{推荐理由：} 不用任何器材就能验证的科学——孩子下次雷雨天自己数秒算距离，成就感极高。

\item[\textbf{第21课}] \textbf{冰为什么能浮在水面上？——水的"反常"性格} \\
\tagphysics \\
\textbf{内容：} 绝大多数物质都是热胀冷缩，水却有"反常"——4℃以上热胀冷缩正常，但从4℃往下降到0℃结冰时反而膨胀，体积变大、密度变小，所以冰浮在水面。这件事对地球生命极其重要：冬天湖面结冰，冰层像一床被子盖在水面上，隔开了冷空气，湖底保持在4℃左右，水里的鱼虾才能活过冬天。如果冰沉底，湖会从底冻到顶，生命全完。延伸到冬天水管为什么会被冰撑爆。\\
\textbf{推荐理由：} 冰块漂在杯子里是日常场景，但这个"反常"的后果极其深刻。理解之后，每次冬天看见冰都会有不同的感受。

\item[\textbf{第22课}] \textbf{指南针为什么指向北？——地球是块大磁铁} \done \\
\tagphysics \\
\textbf{内容：} 战国司南→宋代指南针→郑和与哥伦布。地球内部的铁镍核产生全球磁场；地磁南极在北极附近（反之亦然），指南针N极被地磁S极吸引。磁场形成的"防护罩"挡住太阳风，太阳风在两极形成极光。延伸到候鸟和海龟的"生物指南针"。\\
\textbf{推荐理由：} 指南针是中国古代四大发明之一，有文化自豪感。磁力是另一种"看不见的力"，可以和引力形成呼应。

\item[\textbf{第23课}] \textbf{富兰克林与风筝——电的"真面目"} \done \\
\tagphysics \\
\textbf{内容：} 富兰克林雷雨天放风筝，用钥匙上的火花证明闪电就是静电。电的本质是电子的定向流动。雷电的能量规模。避雷针的工作原理。\textbf{安全用电教育：}不湿手碰电器、不塞金属进插座、远离破损电线、雷雨天不在户外逗留、不超负荷用电、远离高压设施——六条安全规则和"电走它的路，你别搭桥"的总原则。\\
\textbf{推荐理由：} 故事情节紧张刺激，天然吸引注意力。电是孩子最常接触但最不明白的能量形式。安全用电部分将科学认知转化为生存技能，是整模块最实用的内容之一。

\end{enumerate}

\section*{\textcolor{orange!80!black}{五、地球科学与天文}}

\begin{enumerate}

\item[\textbf{第24课}] \textbf{恐龙是怎么灭绝的？——陨石撞击的故事} \\
\tagearth \\
\textbf{内容：} 6600万年前，一颗直径约10公里的小行星撞上今日墨西哥的尤卡坦半岛。撞击释放的能量相当于几十亿颗原子弹，引发全球大火、遮天蔽日的尘埃云（撞击冬天）、酸雨、食物链崩溃。恐龙和其他约75\%的物种灭绝。哺乳动物趁机崛起，最终进化出了人类。\\
\textbf{推荐理由：} 恐龙+陨石是几乎所有孩子都痴迷的话题。灭绝事件是地球科学中最有戏剧性的故事，顺势可引出地质年代、化石、物种灭绝等概念。

\item[\textbf{第25课}] \textbf{天为什么是蓝的？晚霞为什么是红的？——光与大气的"捉迷藏"} \\
\tagearth \\
\textbf{内容：} 太阳光是白光，进入大气后被空气分子向四面八方弹开（瑞利散射）。蓝光波长短，被散射得最厉害——所以天空一片蓝。傍晚太阳角度低、光穿过更厚的大气层，蓝光几乎全被散射掉了，只剩下红光和橙光——所以晚霞是红的。延伸到"星星为什么会眨眼睛"——星光是点光源，穿过动荡的大气层时被不停折射，看起来一闪一闪（而行星是面光源，一般不闪）。\\
\textbf{推荐理由：} 这是孩子每天抬头就能验证的现象。把"天蓝""晚霞""星星眨眼"三个问题串在"光+大气层"一条线上，高效且逻辑清晰。

\item[\textbf{第26课}] \textbf{雨是从哪儿来的？——水循环的故事} \\
\tagearth \\
\textbf{内容：} 一滴水在海里被太阳晒蒸发→变成水蒸气升空→在高空遇冷凝结成小水滴→小水滴聚成云→云飘到陆地上空→水滴变大落下来变成雨→雨水流入河流→河流奔回大海。这就是水循环。延伸到为什么迎风坡降雨多、北京的暴雨是怎么形成的。\\
\textbf{延伸——为什么要节约用水？} 地球表面71\%是水，但能喝的淡水只占不到1\%。水循环让水"再利用"，但人类消耗淡水的速度远快于自然补给的节奏。华北地区人均水资源仅为世界平均水平的1/20。每一滴自来水都要经过取水—净化—输送—加压，背后是能源和资源的消耗。节水不只是"省钱"，是对整个水循环系统的尊重。\\
\textbf{推荐理由：} 下雨天孩子会问"雨从哪儿来"，这是最经典的自然科学问题。节水意识从理解"淡水有多稀缺"开始，而不是从说教开始。

\item[\textbf{第27课}] \textbf{北京为什么春天刮大风沙？——风、沙与防护林} \\
\tagearth \\
\textbf{内容：} 每年3-4月，北京的孩子都会经历沙尘暴。沙子从哪儿来的？（蒙古高原和内蒙古的沙漠/戈壁）。为什么春天刮？（冷暖气团交锋、大风天气）。沙子怎么飞到北京的？国家怎么治理？（三北防护林——"绿色长城"）。延伸到PM10和PM2.5的区别。\\
\textbf{延伸——生态系统观念：} 沙尘暴是一条完整的因果链：过度放牧和开垦→草原退化→土地沙漠化→植被消失→土壤失去固定→风一来沙就飞→沙尘暴。树和草的根就像无数只小手紧紧抓住泥土。植树造林不只是"绿化"，是在修复一个断裂的生态链条。\\
\textbf{推荐理由：} 北京孩子的亲身经历。从"脏脏的天气"引出生态系统观——自然界的每一环都是扣在一起的。

\item[\textbf{第28课}] \textbf{地震是怎么回事？——地球在"活动筋骨"} \\
\tagearth \\
\textbf{内容：} 地球的"外壳"分成好几块板块，像浮在岩浆上的巨大拼图。板块每年移动几厘米，彼此碰撞、挤压、错动——这就是地震的来源。延伸到为什么四川和日本地震多（在板块边界上）、北京也在华北地震带上。\\
\textbf{延伸——地震来了怎么办？} \textbf{伏地、遮挡、抓牢}——躲在结实的桌子下面，护住头部和颈部，远离窗户和吊灯；在室外跑到开阔地带，远离建筑和电线杆；绝不乘坐电梯。破除"站在门框下最安全"的误区（现代建筑门框并不坚固）。\\
\textbf{推荐理由：} 学校有地震演习，孩子需要既知道"为什么地震"又知道"地震了怎么办"。从科学原理到生存技能，一课双收。

\end{enumerate}

\newpage

\section*{\textcolor{blue!70!black}{学科分布统计}}

\begin{center}
\begin{tabular}{|l|c|c|c|}
\hline
\textbf{学科方向} & \textbf{课号} & \textbf{课数} & \textbf{占比} \\
\hline
物理（力与运动） & 第1-3课 & 3 & 11\% \\
生物 & 第4-10课 & 7 & 25\% \\
化学 & 第11-15课 & 5 & 18\% \\
物理（光、热、声、电）& 第16-23课 & 8 & 28\% \\
地球科学与天文 & 第24-28课 & 5 & 18\% \\
\hline
\textbf{合计} & & \textbf{28课} & \\
\hline
\end{tabular}
\end{center}

\vspace{0.3cm}

\textcolor{gray}{\small 注：第11课（肥皂）兼具化学与物理特征，归入化学类。第21课（冰浮水面）兼具物理与化学特征，归入物理类（光热电板块）。第25课涵盖了"星星眨眼"内容。已写7课（第1-4、11、22-23课），待写21课。}

\vspace{0.5cm}

\begin{center}
\rule{0.5\textwidth}{0.5pt}
\vspace{0.3cm}
\textcolor{gray}{\small 大纲版本：2026-05-22 · 待审核}
\end{center}

\end{document}
